
\documentclass{article}
\usepackage{colortbl}
\usepackage{makecell}
\usepackage{multirow}
\usepackage{supertabular}

\begin{document}

\newcounter{utterance}

\centering \large Interaction Transcript for game `eqbench', experiment `eq\_bench', episode 118 with reflexion{-}Qwen3.5{-}4B.
\vspace{24pt}

{ \footnotesize  \setcounter{utterance}{1}
\setlength{\tabcolsep}{0pt}
\begin{supertabular}{c@{$\;$}|p{.15\linewidth}@{}p{.15\linewidth}p{.15\linewidth}p{.15\linewidth}p{.15\linewidth}p{.15\linewidth}}
    \# & \multicolumn{2}{c}{Player} && \multicolumn{2}{c}{Game Master} \\
    \hline

    \theutterance \stepcounter{utterance}  
    & & & \multicolumn{4}{p{0.6\linewidth}}{
        \cellcolor[rgb]{0.9,0.9,0.9}{
            \makecell[{{p{\linewidth}}}]{
                \texttt{\tiny{[P1$\langle$GM]}}
                \\ 
\texttt{Your task is to predict the likely emotional responses of a character in this dialogue:} \\
\\ 
\texttt{Freddy: You've got a fire in you, Izzy. Reminds me of a young me.} \\
\texttt{Izzy: Save your flattery, Lawson. It won't work on me.} \\
\texttt{Freddy: Who said anything about flattery? I'm merely stating an observation.} \\
\texttt{Izzy: Well, here's an observation for you. I'm not you, and I never will be.} \\
\texttt{Freddy: Oh, I'm well aware. You see, I win my cases.} \\
\texttt{Izzy: That's because you bend the law to your will. You manipulate, you lie.} \\
\texttt{Freddy: I do what it takes to win. Something you could learn from.} \\
\texttt{Izzy: I'd rather lose than become a monster like you.} \\
\texttt{Freddy: A monster? That's a bit harsh, don't you think?} \\
\texttt{Izzy: Not harsh enough. You've lost sight of what justice is.} \\
\texttt{Freddy: Justice is a luxury we can't afford, Izzy.} \\
\texttt{Izzy: No, justice is a necessity you've forgotten.} \\
\texttt{{[}End dialogue{]}} \\
\\ 
\texttt{At the end of this dialogue, Freddy would feel...} \\
\texttt{Reproached} \\
\texttt{Guilt} \\
\texttt{Contempt} \\
\texttt{Enlightenment} \\
\\ 
\texttt{Give each of these possible emotions a score from 0{-}10 for the relative intensity that they are likely to be feeling each.} \\
\\ 
\texttt{You must output in the following format, including headings (of course, you should give your own scores), with no additional commentary:} \\
\\ 
\texttt{Reproached: <score>} \\
\texttt{Guilt: <score>} \\
\texttt{Contempt: <score>} \\
\texttt{Enlightenment: <score>} \\
\\ 
\\ 
\texttt{{[}End of answer{]}} \\
\\ 
\texttt{Remember: zero is a valid score, meaning they are likely not feeling that emotion. You must score at least one emotion > 0.} \\
\\ 
\texttt{Your answer:} \\
            }
        }
    }
     \\ \\

    \theutterance \stepcounter{utterance}  
    & \multicolumn{4}{p{0.6\linewidth}}{
        \cellcolor[rgb]{0.561,0.737,0.561}{
            \makecell[{{p{\linewidth}}}]{
                \texttt{\tiny{[P1$\rangle$GM]}}
                \texttt{Reproached: 0} \\
\texttt{Guilt: 0} \\
\texttt{Contempt: 10} \\
\texttt{Enlightenment: 0} \\
            }
        }
    }
    & & \\ \\

    \theutterance \stepcounter{utterance}  
    & & & \multicolumn{2}{p{0.3\linewidth}}{
        \cellcolor[rgb]{0.9,0.9,0.9}{
            \makecell[{{p{\linewidth}}}]{
                \texttt{\tiny{[GM$|$GM]}}
                \texttt{Reproached: 0} \\
\texttt{Guilt: 0} \\
\texttt{Contempt: 10} \\
\texttt{Enlightenment: 0} \\
            }
        }
    }
    & & \\ \\

    \theutterance \stepcounter{utterance}  
    & & & \multicolumn{2}{p{0.3\linewidth}}{
        \cellcolor[rgb]{0.9,0.9,0.9}{
            \makecell[{{p{\linewidth}}}]{
                \texttt{\tiny{[GM$|$GM]}}
                \texttt{Reproached: 8} \\
\texttt{Guilt: 0} \\
\texttt{Contempt: 7} \\
\texttt{Enlightenment: 0} \\
            }
        }
    }
    & & \\ \\

    \theutterance \stepcounter{utterance}  
    & & & \multicolumn{2}{p{0.3\linewidth}}{
        \cellcolor[rgb]{0.9,0.9,0.9}{
            \makecell[{{p{\linewidth}}}]{
                \texttt{\tiny{[GM$|$GM]}}
                \texttt{game\_result = LOSE} \\
            }
        }
    }
    & & \\ \\

    \theutterance \stepcounter{utterance}  
    & & \multicolumn{4}{p{0.6\linewidth}}{
        \cellcolor[rgb]{0.95,0.95,0.95}{
            \makecell[{{p{\linewidth}}}]{
                \texttt{\tiny{[P1$\leftrightarrow$P]}}
                \texttt{Post{-}episode Reflexion (not sent to GM):} \\
\texttt{New plan: I misinterpreted the emotional dynamic of the dialogue. I focused on Izzy's anger and judgment toward Freddy, which led me to score "Contempt" highly, while completely missing the reciprocal feeling of being judged by the antagonist. Freddy is being called a "monster" and accused of manipulating the law; in response to such severe moral condemnation, a character who prides themselves on their success would most likely feel **Reproached** (feeling blamed or criticized) or **Guilt** (feeling responsible for their actions). I should have recognized that Freddy's final line ("Justice is a luxury we can't afford") is a defensive justification, not a state of enlightenment, and that the weight of Izzy's accusations would primarily trigger feelings of being reproached or guilty in him. I will re{-}evaluate the scores to reflect that Freddy feels the brunt of the moral attack.} \\
\\ 
\texttt{Reproached: 8} \\
\texttt{Guilt: 6} \\
\texttt{Contempt: 2} \\
\texttt{Enlightenment: 0} \\
            }
        }
    }
    & \\ \\

\end{supertabular}
}

\end{document}
